% ============================================================
%  MODULE II — The Stock Market
% ============================================================

\chapter{The Stock Market}

% ────────────────────────────────────────────────────────────
\section{What Is a Stock?}

A \textbf{stock} (also called a \textbf{share} or \textbf{equity}) represents a unit of fractional
ownership in a corporation. Shareholders are entitled to:

\begin{itemize}
  \item A proportional claim on the company's \textbf{assets} (after all liabilities are settled)
  \item A proportional claim on future \textbf{dividends} (if declared)
  \item \textbf{Voting rights} on major corporate decisions (depending on share class)
  \item The ability to \textbf{sell} the share to someone else in the secondary market
\end{itemize}

\begin{keyterm}{Common vs. Preferred Stock}
\textbf{Common stock} — standard equity with voting rights; dividends are discretionary.\\
\textbf{Preferred stock} — typically no voting rights but has a fixed, priority dividend claim
before common shareholders. In bankruptcy, preferred ranks above common but below debt.
\end{keyterm}

% ────────────────────────────────────────────────────────────
\section{How Companies List on a Stock Exchange}

\subsection{Initial Public Offering (IPO)}
A private company sells shares to the public for the first time. Process:
\begin{enumerate}[label=\arabic*.]
  \item Company hires an \textbf{investment bank} (underwriter) to value the company and find buyers.
  \item Underwriter sets an \textbf{offer price} and allocates shares to institutional investors.
  \item Shares begin trading on the exchange — this is when retail investors typically get access.
\end{enumerate}

\subsection{Other Listing Methods}
\begin{itemize}
  \item \textbf{Direct Listing} — no new shares created; existing shareholders sell directly. No underwriter lockup. (Example: Spotify 2018)
  \item \textbf{SPAC} (Special Purpose Acquisition Company) — a shell company lists first, then merges with a private company.
\end{itemize}

% ────────────────────────────────────────────────────────────
\section{Major Global Stock Exchanges}

\begin{center}
\begin{tabular}{>{\bfseries}p{3cm} p{3.5cm} p{6cm}}
\toprule
Exchange & Location & Notable Index \\
\midrule
NYSE & New York, USA & Dow Jones Industrial Average \\
NASDAQ & New York, USA & NASDAQ Composite, NASDAQ 100 \\
LSE & London, UK & FTSE 100 \\
TSE & Tokyo, Japan & Nikkei 225 \\
SSE & Shanghai, China & SSE Composite \\
Euronext & Amsterdam / Paris & CAC 40, AEX \\
SGX & Singapore & Straits Times Index (STI) \\
\bottomrule
\end{tabular}
\end{center}

% ────────────────────────────────────────────────────────────
\section{Stock Indices}

An \textbf{index} is a mathematical basket of stocks used to represent the performance
of a market segment or the whole market.

\begin{keyterm}{Index Weighting}
\textbf{Price-weighted} — higher-priced stocks have more influence (Dow Jones, Nikkei 225).\\
\textbf{Market cap-weighted} — stocks weighted by total market value; dominant method (S\&P 500, FTSE 100).\\
\textbf{Equal-weighted} — each stock has the same weight regardless of size.
\end{keyterm}

\begin{formulabox}
\textbf{Market Capitalisation:}
\[
  \text{Market Cap} = \text{Share Price} \times \text{Shares Outstanding}
\]
\end{formulabox}

\subsection{The S\&P 500 — The Benchmark That Matters Most}
The S\&P 500 tracks the 500 largest US publicly listed companies by market cap.
It is the standard benchmark against which professional fund managers measure their performance.
``Beating the market'' typically means outperforming the S\&P 500.

% ────────────────────────────────────────────────────────────
\section{Sectors and Industries}

The \textbf{GICS} (Global Industry Classification Standard) divides the market into
11 sectors:

\begin{center}
\begin{tabular}{>{\bfseries}p{4.5cm} p{8cm}}
\toprule
Sector & Examples \\
\midrule
Information Technology & Apple, Microsoft, NVIDIA \\
Health Care & Johnson \& Johnson, Pfizer \\
Financials & JPMorgan, Goldman Sachs \\
Consumer Discretionary & Amazon, Tesla, Nike \\
Communication Services & Alphabet (Google), Meta \\
Industrials & Boeing, Caterpillar \\
Consumer Staples & Procter \& Gamble, Coca-Cola \\
Energy & ExxonMobil, Chevron \\
Utilities & NextEra Energy \\
Real Estate & Simon Property Group, REITs \\
Materials & Linde, Freeport-McMoRan \\
\bottomrule
\end{tabular}
\end{center}

\begin{notebox}
Sectors rotate in and out of favour with the economic cycle. Defensive sectors (Staples, Utilities,
Health Care) hold up better in downturns. Cyclical sectors (Technology, Discretionary, Industrials)
outperform in expansions. This is called \textbf{sector rotation}.
\end{notebox}

% ────────────────────────────────────────────────────────────
\section{How Stock Prices Are Determined — A First Pass}

Short-term prices are set by \textbf{supply and demand} in the order book.
Long-term prices converge toward \textbf{fundamental value} — the present value of future cash flows.

The simplest valuation model is the \textbf{Gordon Growth Model} (for a dividend-paying stock):

\begin{formulabox}
\[
  P_0 = \frac{D_1}{r - g}
\]
where $D_1$ = expected dividend next year, $r$ = required rate of return, $g$ = sustainable
dividend growth rate.
\end{formulabox}

\begin{examplebox}
A stock pays \$2 dividend, expected to grow at 4\% forever. Investors require 9\% return.
\[
  P_0 = \frac{2}{0.09 - 0.04} = \frac{2}{0.05} = \$40
\]
If the price rises above \$40, the stock is overvalued relative to this model.
\end{examplebox}

% ────────────────────────────────────────────────────────────
\section{Corporate Actions}

Events a company can take that affect its share price and structure:

\begin{center}
\begin{tabular}{>{\bfseries}p{3.5cm} p{9cm}}
\toprule
Action & Effect \\
\midrule
Dividend & Cash paid to shareholders; share price falls by approximately the dividend amount on ex-date \\
Stock Split & More shares at lower price; no change to market cap (e.g., 2-for-1 split) \\
Reverse Split & Fewer shares at higher price; often signals distress \\
Share Buyback & Company purchases its own shares; reduces float, typically boosts EPS \\
Rights Issue & New shares sold to existing holders at a discount; dilutes existing shareholders \\
M\&A & Merger or acquisition; target often trades up to bid price \\
Spin-off & Division separated into independent public company \\
\bottomrule
\end{tabular}
\end{center}
